% ============================================================================
% ANHANG M: Die inhärente Metrik der Weber-Elektrodynamik und Weber-Gravitation
% ============================================================================

\chapter{Die inhärente Metrik der Weber-Elektrodynamik und Weber-Gravitation}
\label{app:weber_metrik}

\section{Einleitung}
\label{app:weber_metrik_einleitung}

Die Weber-Elektrodynamik (WED) und die Weber-Gravitation (WG) besitzen eine fundamentale Eigenschaft, die sie von der Maxwell-Theorie und der Allgemeinen Relativitätstheorie (ART) unterscheidet: \textbf{Die Weber'sche Lagrange-Funktion enthält eine inhärente, geschwindigkeitsunabhängige Metrik, die am kritischen Radius singular wird und einen Übergang von Riemannscher zu Lorentzscher Geometrie beschreibt.} \cite{FrauenfelderWeber2024}

Diese Eigenschaft zeigt, dass die WED und die WG nicht nur alternative Kraftbeschreibungen sind, sondern vollständige geometrische Formalismen, die ihr eigenes Raummodell in der mathematischen Struktur tragen. Die folgende Darstellung folgt der Arbeit von Frauenfelder und Weber \cite{FrauenfelderWeber2024} und ordnet sie in den Rahmen der IWT/WDBT+ ein.

\section{Die Weber'sche Lagrange-Funktion der WED}
\label{app:weber_metrik_lagrange_wed}

Die Weber'sche Lagrange-Funktion für zwei Ladungen in Polarkoordinaten lautet (siehe auch Kapitel~\ref{ch:dstt_weber_kraefte}) \cite{Assis2005}:

\[
L_{\mathrm{W}}(r,\phi,v_r,v_\phi) = \frac{1}{2}(v_r^2 + r^2 v_\phi^2) - \frac{1}{r}\left(1 + \frac{v_r^2}{2c^2}\right)
\tag{M.1}
\]

Der erste Term ist die übliche kinetische Energie, der zweite Term beschreibt ein geschwindigkeitsabhängiges Potential. Historisch führte dieser Term zu Verwirrung, insbesondere bei Helmholtz, der die Energieerhaltung anzweifelte \cite{Assis2005}. Die moderne Forschung hat diese Einwände jedoch widerlegt \cite{Assis2012}.

\section{Die Umformulierung und die inhärente Metrik der WED}
\label{app:weber_metrik_umformulierung_wed}

Durch einfaches Ausmultiplizieren und Umordnen der Terme ergibt sich eine völlig neue Interpretation \cite{FrauenfelderWeber2024}:

\[
L_{\mathrm{W}}(r,\phi,v_r,v_\phi) = 
\frac{1}{2}\left(1 - \frac{1}{c^2r}\right) v_r^2 + \frac{1}{2} r^2 v_\phi^2 - \frac{1}{r}
\tag{M.2}
\]

oder kompakt:

\[
L_{\mathrm{W}} = \frac{1}{2}\left( g_{rr}\, v_r^2 + g_{\phi\phi}\, v_\phi^2 \right) - \frac{1}{r}
\tag{M.3}
\]

mit:

\[
g_{rr} = 1 - \frac{1}{c^2r}, \qquad g_{\phi\phi} = r^2
\tag{M.4}
\]

Die entscheidende Konsequenz: \textbf{Das Potential hängt nun nicht mehr von der Geschwindigkeit ab. Die gesamte geschwindigkeitsabhängige Dynamik steckt in der Metrik.} \cite{FrauenfelderWeber2024} Dies ist ein zentrales Ergebnis, das die WED von der Maxwell-Theorie fundamental unterscheidet.

\section{Die Metrik und der kritische Radius der WED}
\label{app:weber_metrik_metrik_wed}

Die Metrik hat die Form \cite{FrauenfelderWeber2024}:

\[
g_{\text{WED}} = \left(1 - \frac{1}{c^2r}\right) dr^2 + r^2 d\phi^2
\tag{M.5}
\]

Der \textbf{kritische Radius} (Weber-Radius) ist definiert als:

\[
r_c = \frac{1}{c^2}
\tag{M.6}
\]

Die physikalische Bedeutung \cite{FrauenfelderWeber2024}:

\begin{table}[htbp]
\centering
\begin{tabular}{|c|c|c|}
\hline
\textbf{Bereich} & \textbf{Metrik} & \textbf{Geometrie} \\
\hline
\( r > r_c \) & \( g_{rr} > 0 \) & \textbf{Riemannsch} (gekrümmter Raum, ART-ähnlich) \\
\hline
\( r = r_c \) & \( g_{rr} = 0 \) & \textbf{Singularität} (Übergang) \\
\hline
\( r < r_c \) & \( g_{rr} < 0 \) & \textbf{Lorentzsch} (Raumzeit-ähnlich, SRT-ähnlich) \\
\hline
\end{tabular}
\caption{Die drei Bereiche der Weber-Metrik der WED.}
\label{tab:weber_metrik_bereiche_wed}
\end{table}

\textbf{Zentrale Aussage:} Die Weber'sche Lagrange-Funktion enthält eine inhärente Metrik, die am kritischen Radius \( r_c \) von Riemannscher zu Lorentzscher Geometrie wechselt \cite{FrauenfelderWeber2024}. Diese Metrik ist nicht willkürlich hinzugefügt, sondern eine direkte Konsequenz der mathematischen Struktur der WED.

\section{Die inhärente Metrik der Weber-Gravitation (WG)}
\label{app:weber_metrik_wg}

Die Weber-Gravitation (WG) für zwei Massen \( M \) und \( m \) hat eine analoge Struktur (siehe Kapitel~\ref{ch:dstt_weber_kraefte}) \cite{Assis2005}:

\[
\vec{F}_{\text{WG}} = -\frac{GMm}{r^2}
\left\{
\left[
1 - \frac{\dot{r}^2}{c^2} + \frac{1}{2}\frac{r\ddot{r}}{c^2}
\right] \hat{r}
+ \frac{2(\dot{r} \cdot \vec{v})}{c^2} \vec{v}
\right\}
\tag{M.7}
\]

Die Lagrange-Funktion der WG lautet \cite{Assis2005}:

\[
L_{\text{WG}}(r,\phi,v_r,v_\phi) = \frac{1}{2}(v_r^2 + r^2 v_\phi^2) + \frac{GM}{r}
\left(1 - \frac{v_r^2}{2c^2}\right)
\tag{M.8}
\]

Durch dieselbe Umformulierung wie bei der WED ergibt sich \cite{FrauenfelderWeber2024}:

\[
L_{\text{WG}} = \frac{1}{2}\left(1 - \frac{GM}{c^2 r}\right) v_r^2 + \frac{1}{2} r^2 v_\phi^2 + \frac{GM}{r}
\tag{M.9}
\]

oder kompakt:

\[
L_{\text{WG}} = \frac{1}{2}\left( g_{rr}\, v_r^2 + g_{\phi\phi}\, v_\phi^2 \right) + \frac{GM}{r}
\tag{M.10}
\]

mit:

\[
g_{rr} = 1 - \frac{GM}{c^2 r}, \qquad g_{\phi\phi} = r^2
\tag{M.11}
\]

Die Metrik der WG hat die Form \cite{FrauenfelderWeber2024}:

\[
g_{\text{WG}} = \left(1 - \frac{GM}{c^2 r}\right) dr^2 + r^2 d\phi^2
\tag{M.12}
\]

Der \textbf{kritische Radius} der WG (Weber-Schwarzschild-Radius) ist definiert als:

\[
r_{s,\text{Weber}} = \frac{GM}{c^2}
\tag{M.13}
\]

Die physikalische Bedeutung \cite{FrauenfelderWeber2024}:

\begin{table}[htbp]
\centering
\begin{tabular}{|c|c|c|}
\hline
\textbf{Bereich} & \textbf{Metrik} & \textbf{Geometrie} \\
\hline
\( r > r_{s,\text{Weber}} \) & \( g_{rr} > 0 \) & \textbf{Riemannsch} (gekrümmter Raum, ART-ähnlich) \\
\hline
\( r = r_{s,\text{Weber}} \) & \( g_{rr} = 0 \) & \textbf{Singularität} (Übergang) \\
\hline
\( r < r_{s,\text{Weber}} \) & \( g_{rr} < 0 \) & \textbf{Lorentzsch} (Raumzeit-ähnlich, SRT-ähnlich) \\
\hline
\end{tabular}
\caption{Die drei Bereiche der Weber-Metrik der WG.}
\label{tab:weber_metrik_bereiche_wg}
\end{table}

\textbf{Zentrale Aussage:} Die Weber-Gravitation enthält ebenfalls eine inhärente Metrik, die am kritischen Radius \( r_{s,\text{Weber}} = GM/c^2 \) von Riemannscher zu Lorentzscher Geometrie wechselt \cite{FrauenfelderWeber2024}. Dieser kritische Radius entspricht formal dem Schwarzschild-Radius der ART, hat aber eine völlig andere physikalische Interpretation: Er ist nicht der Ereignishorizont eines Schwarzen Lochs, sondern der Übergangspunkt zwischen zwei geometrischen Bereichen derselben Theorie.

\section{Physikalische Konsequenzen für ART und SRT}
\label{app:weber_metrik_konsequenzen}

Aus dieser inhärenten Metrik der WED und WG folgen vier zentrale physikalische Konsequenzen \cite{FrauenfelderWeber2024,Assis2005,Assis2012}:

\begin{enumerate}
    \item \textbf{Die WED und WG enthalten ihr eigenes Raummodell:} Sie benötigen keine äußere Raumzeit, sondern tragen die Geometrie in sich selbst. Dies bedeutet: Die Maxwell-Theorie und die ART sind \textbf{nicht fundamental}, sondern emergente Grenzfälle der WED und WG.
    
    \item \textbf{ART und SRT sind in der WED/WG vereint:} Die Metrik beschreibt beide Bereiche – den gekrümmten Raum (Riemannsch, \( r > r_c \)) und die Lorentz'sche Raumzeit (\( r < r_c \)) – als zwei Seiten derselben mathematischen Struktur. \textbf{Die ART und die SRT sind damit nicht zwei separate Theorien, sondern zwei verschiedene geometrische Bereiche einer einzigen Theorie.} Dies bedeutet:
    
    \begin{itemize}
        \item Die \textbf{ART} ist der Grenzfall \( r > r_s \): gekrümmter Raum, Gravitation als Geometrie.
        \item Die \textbf{SRT} ist der Grenzfall \( r < r_s \): flache Raumzeit, Lorentz-Transformationen.
        \item Die \textbf{WG} ist die vollständige Theorie, die beide Bereiche vereint.
    \end{itemize}
    
    Dies ist ein fundamentaler Unterschied zur Maxwell-Theorie, die keine solche Metrik enthält, und zur ART, die eine gekrümmte Raumzeit postuliert, ohne sie aus einer tieferen Struktur abzuleiten.
    
    \item \textbf{Keine geschlossenen Bahnen im Inneren:} Innerhalb des kritischen Radius gibt es keine periodischen Bahnen; die Trajektorien spiralen in den Ursprung. Dies bedeutet, dass es in der WG keine stabilen inneren Bahnen gibt – eine wichtige Abweichung von der ART, die geschlossene Bahnen innerhalb des Schwarzschild-Radius (für \( r < 3r_s \)) zulässt.
    
    \item \textbf{Quantisierung möglich:} Die Lorentz'sche Interpretation öffnet die Tür zur Schrödinger-Gleichung für den Weber-Kern, indem die kinetische Energie durch den Laplace-Beltrami-Operator ersetzt wird \cite{FrauenfelderWeber2024}. Dies ist ein Hinweis darauf, dass die WED und WG eine natürliche Brücke zur Quantenmechanik bilden.
\end{enumerate}

\section{Die Implikationen für die Maxwell-Theorie und die ART}
\label{app:weber_metrik_maxwell_art}

Aus der Existenz der inhärenten Metrik in der WED und WG folgt zwingend:

\begin{enumerate}
    \item \textbf{Die Maxwell-Theorie ist ein Grenzfall der WED:} Die Maxwell-Theorie postuliert Felder \( \vec{E} \) und \( \vec{B} \) als fundamentale Größen. Die WED zeigt, dass diese Felder emergente Phänomene sind, die aus der direkten Wechselwirkung von Ladungen folgen. Die Maxwell-Theorie ist daher \textbf{nicht fundamental}.
    
    \item \textbf{Die ART ist ein Grenzfall der WG:} Die ART postuliert eine gekrümmte Raumzeit als fundamental. Die WG zeigt, dass diese Krümmung aus der inhärenten Metrik der direkten Massenwechselwirkung folgt. Die ART ist daher \textbf{nicht fundamental}.
    
    \item \textbf{Die ART und die SRT sind in der WG vereint:} Die WG enthält sowohl den Riemannschen Bereich (\( r > r_s \)) als auch den Lorentzschen Bereich (\( r < r_s \)). Die ART und die SRT sind damit keine zwei separaten Theorien, sondern zwei verschiedene geometrische Bereiche einer einzigen Theorie.
    
    \item \textbf{Die WED und die WG sind fundamental:} Sie enthalten ihr eigenes Raummodell, benötigen keine äußere Raumzeit und sind mathematisch konsistent. Sie sind daher die fundamentaleren Theorien.
\end{enumerate}

\section{Verbindung zur IWT/WDBT+}
\label{app:weber_metrik_wdbt}

Die hier beschriebene inhärente Metrik der WED und WG ist ein zentraler Baustein der Informations-Weber-Theorie (IWT) und der Weber-de-Broglie-Bohm-Theorie mit fraktalem Raum (WDBT+). 

In der IWT emergiert die Metrik aus der fraktalen Kopplungsmatrix \( K_{kl}^{(n)} = 1/d_{kl}^{3-D} \) (siehe Kapitel~\ref{ch:emergenz_raum_zeit} und Anhang~\ref{app:mathematik}). Die kontinuierliche Näherung dieser diskreten Metrik führt im Grenzfall \( D = 2 \) auf die hier beschriebene Weber-Metrik der WED und WG.

Die fraktale Raumdimension \( D \approx 2,704 \) (siehe Kapitel~\ref{ch:fraktale_dimension}) ist damit eine Verallgemeinerung der Weber-Metrik auf den fraktalen Raum. Die Korrelationslänge \( L_{Q,0} \approx 2{,}0 \times 10^{46} \, \text{m} \) (siehe Anhang~\ref{app:neutrino_oszillation}) ist die maximale physikalische Skala dieser Metrik.

\section{Zusammenfassung}
\label{app:weber_metrik_zusammenfassung}

Die Weber-Elektrodynamik (WED) und die Weber-Gravitation (WG) enthalten inhärente Metriken:

\[
g_{\text{WED}} = \left(1 - \frac{1}{c^2r}\right) dr^2 + r^2 d\phi^2
\]

\[
g_{\text{WG}} = \left(1 - \frac{GM}{c^2 r}\right) dr^2 + r^2 d\phi^2
\]

Diese Metriken:
\begin{itemize}
    \item sind eine direkte Konsequenz der mathematischen Struktur der WED und WG,
    \item werden am kritischen Radius \( r_c = 1/c^2 \) (WED) bzw. \( r_s = GM/c^2 \) (WG) singular,
    \item wechseln von Riemannscher zu Lorentzscher Geometrie,
    \item sind nicht willkürlich hinzugefügt, sondern inhärent,
    \item ermöglichen eine konsistente Quantisierung des Weber-Kerns,
    \item sind der kontinuierliche Grenzfall der diskreten Metrik der IWT,
    \item werden in der WDBT+ auf den fraktalen Raum verallgemeinert.
\end{itemize}

\textbf{Die zentrale Aussage für ART und SRT:}

Die ART und die SRT sind \textbf{nicht fundamental}. Sie sind zwei verschiedene geometrische Bereiche einer einzigen, tieferen Theorie – der Weber-Gravitation. Die WG enthält sowohl den Riemannschen Bereich (\( r > r_s \)) als auch den Lorentzschen Bereich (\( r < r_s \)). Damit sind die ART und die SRT in der WG vereint – nicht als separate Theorien, sondern als zwei Aspekte derselben mathematischen Struktur.

Damit besitzen die WED und WG eigene Raummodelle, die sie von der Maxwell-Theorie und der ART grundlegend unterscheiden. Die Maxwell-Theorie und die ART sind emergente Grenzfälle der WED und WG – nicht umgekehrt.